\documentclass[aspectratio=169,11pt]{beamer}
\usetheme{Madrid}
\usecolortheme{dolphin}
\setbeamertemplate{navigation symbols}{}
\setbeamertemplate{caption}[numbered]

\usepackage[utf8]{inputenc}
\usepackage[T1]{fontenc}
\usepackage[spanish,es-noshorthands]{babel}
\usepackage{amsmath,amssymb,amsthm,mathtools}
\usepackage{tikz}
\usepackage{pgfplots}
\pgfplotsset{compat=1.18}
\usepackage{booktabs}
\usepackage{array}

\title[Prob. y Estadística]{Clase 30: Inferencia en el modelo de regresión}
\subtitle{Unidad 8: Modelos lineales}
\institute[UNS]{Universidad Nacional del Sur \\ Departamento de Matemática}
\author{Probabilidad y Estadística (5765)}
\date{}

\begin{document}

\begin{frame}
\titlepage
\end{frame}

\begin{frame}{Contenidos}
\tableofcontents
\end{frame}

\section{El supuesto de normalidad}

\begin{frame}{De la estimación puntual a la inferencia}
En la Clase 29 obtuvimos los estimadores de mínimos cuadrados $\hat\beta_0,\hat\beta_1$ y vimos que son insesgados. Para construir \textbf{pruebas de hipótesis} e \textbf{intervalos de confianza} necesitamos conocer su \textbf{distribución muestral}, lo cual requiere un supuesto adicional.

\begin{block}{Supuesto de normalidad}
\[
\varepsilon_i \sim N(0,\sigma^2), \qquad \varepsilon_1,\dots,\varepsilon_n \text{ independientes},
\]
o, equivalentemente,
\[
Y_i \sim N(\beta_0+\beta_1 x_i,\ \sigma^2), \qquad Y_1,\dots,Y_n \text{ independientes}.
\]
\end{block}

Bajo este supuesto, como $\hat\beta_0$ y $\hat\beta_1$ son \textbf{combinaciones lineales} de las $Y_i$, también tienen distribución normal.
\end{frame}

\section{Distribución de los estimadores}

\begin{frame}{Distribución de $\hat\beta_1$}
Recordando que $\hat\beta_1 = \sum_i c_i Y_i$ con $c_i=(x_i-\bar x)/S_{xx}$, y que una combinación lineal de normales independientes es normal:

\begin{block}{Distribución de $\hat\beta_1$}
\[
\hat\beta_1 \sim N\!\left(\beta_1,\ \frac{\sigma^2}{S_{xx}}\right).
\]
\end{block}

\textit{Esperanza:} ya vimos $E(\hat\beta_1)=\beta_1$ (Clase 29).

\textit{Varianza:}
\[
V(\hat\beta_1) = \sum_i c_i^2 V(Y_i) = \sigma^2 \sum_i c_i^2 = \sigma^2 \sum_i \frac{(x_i-\bar x)^2}{S_{xx}^2} = \frac{\sigma^2}{S_{xx}}.
\]

A mayor dispersión de los $x_i$ (mayor $S_{xx}$), menor es la varianza de $\hat\beta_1$: los datos más ``separados'' en $X$ estiman la pendiente con mayor precisión.
\end{frame}

\begin{frame}{Distribución de $\hat\beta_0$}
De manera análoga, usando $\hat\beta_0 = \bar Y - \hat\beta_1\bar x$:

\begin{block}{Distribución de $\hat\beta_0$}
\[
\hat\beta_0 \sim N\!\left(\beta_0,\ \sigma^2\left[\frac{1}{n}+\frac{\bar x^2}{S_{xx}}\right]\right).
\]
\end{block}

Ambos estimadores son normales, insesgados, pero \textbf{no independientes} entre sí (están correlacionados a través de $\bar x$), salvo en el caso particular $\bar x = 0$.
\end{frame}

\section{Estimación de la varianza del error}

\begin{frame}{Cuadrado medio del error}
El parámetro $\sigma^2$ (varianza del error) es desconocido y debe estimarse a partir de los residuos.

\begin{definition}
El \textbf{cuadrado medio del error} es
\[
CME = \frac{SCE}{n-2} = \frac{\sum_{i=1}^n (y_i-\hat y_i)^2}{n-2} = \frac{\sum_{i=1}^n e_i^2}{n-2}.
\]
\end{definition}

\begin{block}{¿Por qué $n-2$ grados de libertad?}
Se ``pierden'' 2 grados de libertad porque $SCE$ se calcula a partir de los residuos de una recta ajustada con \textbf{2} parámetros estimados ($\hat\beta_0$ y $\hat\beta_1$). Puede probarse que
\[
E(CME) = \sigma^2,
\]
es decir, $CME$ es un estimador \textbf{insesgado} de $\sigma^2$. Se define además $\hat\sigma = \sqrt{CME}$, el \textbf{error estándar de la estimación}.
\end{block}
\end{frame}

\begin{frame}{Errores estándar de $\hat\beta_0$ y $\hat\beta_1$}
Reemplazando $\sigma^2$ por su estimador $CME$ obtenemos los \textbf{errores estándar estimados}:
\[
\widehat{ee}(\hat\beta_1) = \sqrt{\frac{CME}{S_{xx}}}, \qquad
\widehat{ee}(\hat\beta_0) = \sqrt{CME\left[\frac{1}{n}+\frac{\bar x^2}{S_{xx}}\right]}.
\]

\begin{block}{Resultado clave para la inferencia}
Bajo normalidad de los errores,
\[
\frac{\hat\beta_1 - \beta_1}{\widehat{ee}(\hat\beta_1)} \sim t_{n-2}, \qquad
\frac{\hat\beta_0 - \beta_0}{\widehat{ee}(\hat\beta_0)} \sim t_{n-2},
\]
distribución $t$ de Student con $n-2$ grados de libertad (los mismos que en la estimación de una media con $\sigma$ desconocida, pero aquí ``se gastan'' 2 grados de libertad en vez de 1).
\end{block}
\end{frame}

\section{Inferencia sobre la pendiente $\beta_1$}

\begin{frame}{Intervalo de confianza para $\beta_1$}
\begin{block}{IC del $(1-\alpha)100\%$ para $\beta_1$}
\[
\hat\beta_1 \pm t_{n-2,\,\alpha/2}\ \widehat{ee}(\hat\beta_1).
\]
\end{block}

\begin{block}{Prueba de hipótesis para $\beta_1$}
La hipótesis de mayor interés práctico es
\[
H_0: \beta_1 = 0 \qquad \text{vs.} \qquad H_1: \beta_1 \neq 0
\]
(u opciones unilaterales), ya que $\beta_1=0$ significa que $X$ \textbf{no aporta información lineal} sobre $Y$ (no hay relación lineal). El estadístico de prueba es
\[
T = \frac{\hat\beta_1 - 0}{\widehat{ee}(\hat\beta_1)} \sim t_{n-2} \quad \text{bajo } H_0,
\]
y se rechaza $H_0$ si $|T| > t_{n-2,\alpha/2}$ (bilateral) al nivel $\alpha$.
\end{block}
\end{frame}

\begin{frame}{Inferencia sobre $\beta_0$}
De manera totalmente análoga:

\begin{block}{IC e hipótesis para $\beta_0$}
\[
\text{IC: } \hat\beta_0 \pm t_{n-2,\,\alpha/2}\ \widehat{ee}(\hat\beta_0), \qquad
T = \frac{\hat\beta_0-\beta_{0,0}}{\widehat{ee}(\hat\beta_0)} \sim t_{n-2} \ \text{bajo } H_0:\beta_0=\beta_{0,0}.
\]
\end{block}

\begin{alertblock}{Observación}
En la mayoría de las aplicaciones el interés central está en $\beta_1$ (existencia y magnitud de la relación lineal), ya que $\beta_0$ suele carecer de interpretación práctica directa (con frecuencia representa una extrapolación fuera del rango de los datos).
\end{alertblock}
\end{frame}

\section{Significancia global de la regresión}

\begin{frame}{Prueba F de significancia de la regresión}
En regresión lineal \textbf{simple} (una sola regresora), probar $H_0:\beta_1=0$ mediante el estadístico $t$ es equivalente a una prueba \textbf{F} basada en la descomposición $SCT=SCR+SCE$:

\[
F = \frac{SCR/1}{SCE/(n-2)} = \frac{CMR}{CME} \sim F_{1,\,n-2} \quad \text{bajo } H_0:\beta_1=0.
\]

\begin{block}{Relación entre $T$ y $F$}
Puede demostrarse que $F = T^2$, de modo que ambas pruebas llevan siempre a la misma conclusión en regresión simple. La prueba $F$ es la que se generaliza naturalmente a regresión \textbf{múltiple} y al análisis de la varianza (Clases 31--32).
\end{block}
\end{frame}

\begin{frame}{Tabla ANOVA de la regresión}
Es habitual resumir la descomposición de la variabilidad en una \textbf{tabla ANOVA de regresión}:

\begin{table}[htbp]
\centering
\begin{tabular}{lcccc}
\toprule
Fuente & SC & gl & CM & F \\
\midrule
Regresión & $SCR$ & $1$ & $CMR = SCR/1$ & $CMR/CME$ \\
Error & $SCE$ & $n-2$ & $CME = SCE/(n-2)$ & \\
\midrule
Total & $SCT$ & $n-1$ & & \\
\bottomrule
\end{tabular}
\caption{Tabla ANOVA para la regresión lineal simple}
\end{table}

Nótese que los grados de libertad también se descomponen: $n-1 = 1 + (n-2)$.
\end{frame}

\section{Predicción}

\begin{frame}{Predicción de la respuesta media}
Para un valor particular $x_0$ de la regresora, distinguimos dos problemas de predicción distintos.

\begin{block}{(a) Respuesta media en $x_0$}
Estimamos $E(Y\mid X=x_0)=\beta_0+\beta_1 x_0$ mediante $\hat y_0 = \hat\beta_0+\hat\beta_1 x_0$, con
\[
V(\hat y_0) = \sigma^2\left[\frac{1}{n}+\frac{(x_0-\bar x)^2}{S_{xx}}\right].
\]
IC del $(1-\alpha)100\%$ para la respuesta media:
\[
\hat y_0 \pm t_{n-2,\alpha/2}\sqrt{CME\left[\frac{1}{n}+\frac{(x_0-\bar x)^2}{S_{xx}}\right]}.
\]
\end{block}
\end{frame}

\begin{frame}{Predicción de una observación individual}
\begin{block}{(b) Valor individual futuro en $x_0$}
Predecimos una \textbf{nueva} observación $Y_0 = \beta_0+\beta_1 x_0+\varepsilon_0$ con el mismo punto $\hat y_0$, pero ahora hay que sumar también la variabilidad propia del nuevo error $\varepsilon_0$:
\[
V(Y_0-\hat y_0) = \sigma^2\left[1+\frac{1}{n}+\frac{(x_0-\bar x)^2}{S_{xx}}\right].
\]
\textbf{Intervalo de predicción} del $(1-\alpha)100\%$:
\[
\hat y_0 \pm t_{n-2,\alpha/2}\sqrt{CME\left[1+\frac{1}{n}+\frac{(x_0-\bar x)^2}{S_{xx}}\right]}.
\]
\end{block}

\begin{alertblock}{Comparación}
El intervalo de \textbf{predicción} (b) es siempre más ancho que el de \textbf{respuesta media} (a), porque incorpora tanto la incertidumbre en la estimación de la recta como la variabilidad individual de una observación nueva. Ambos son más angostos cuanto más cerca está $x_0$ de $\bar x$.
\end{alertblock}
\end{frame}

\section{Ejemplo numérico completo}

\begin{frame}{Ejemplo: retomando la soldadura}
Retomamos los datos de la Clase 29: tiempo de soldadura ($X$, segundos) y resistencia a la tracción ($Y$, kg), $n=8$.

\begin{table}[htbp]
\centering
\begin{tabular}{c|cccccccc}
\toprule
$x_i$ & 2 & 3 & 4 & 5 & 6 & 7 & 8 & 9 \\
$y_i$ & 32 & 38 & 41 & 45 & 48 & 53 & 55 & 61 \\
\bottomrule
\end{tabular}
\end{table}

Ya sabemos: $\hat\beta_0\approx 25.214$, $\hat\beta_1\approx 3.893$, $S_{xx}=42$, $S_{yy}=641.875$, $S_{xy}=163.5$, $\bar x=5.5$.

\[
SCR = \hat\beta_1\, S_{xy} = \frac{S_{xy}^2}{S_{xx}} = \frac{163.5^2}{42} = \frac{26732.25}{42} \approx 636.482.
\]
\[
SCE = SCT - SCR = S_{yy} - SCR = 641.875 - 636.482 \approx 5.393.
\]
\end{frame}

\begin{frame}{Ejemplo: tabla ANOVA de la regresión}
\begin{table}[htbp]
\centering
\begin{tabular}{lcccc}
\toprule
Fuente & SC & gl & CM & F \\
\midrule
Regresión & $636.482$ & $1$ & $636.482$ & $708.1$ \\
Error & $5.393$ & $6$ & $0.899$ & \\
\midrule
Total & $641.875$ & $7$ & & \\
\bottomrule
\end{tabular}
\end{table}

\[
CME = \frac{SCE}{n-2} = \frac{5.393}{6} \approx 0.899, \qquad F = \frac{CMR}{CME} = \frac{636.482}{0.899} \approx 708.1.
\]

Con $F_{1,6;0.05}=5.99$: como $F=708.1 \gg 5.99$, se \textbf{rechaza} $H_0:\beta_1=0$: la regresión es altamente significativa.
\end{frame}

\begin{frame}{Ejemplo: IC para $\beta_1$ y predicción}
Con $CME=0.899$, $\widehat{ee}(\hat\beta_1)=\sqrt{0.899/42}=\sqrt{0.02140}\approx0.1463$.

IC del 95\% para $\beta_1$ ($t_{6;0.025}=2.447$):
\[
3.893 \pm 2.447(0.1463) = 3.893\pm0.358 \ \Rightarrow\ (3.535,\ 4.251).
\]
Como el intervalo no contiene al $0$, se confirma (por otra vía) el rechazo de $H_0:\beta_1=0$.

\textbf{Predicción en $x_0=6.5$} (respuesta media): $\hat y_0 = 25.214+3.893(6.5)\approx50.518$ kg, con
\[
\widehat{ee}(\hat y_0)=\sqrt{0.899\left[\tfrac18+\tfrac{(6.5-5.5)^2}{42}\right]} = \sqrt{0.899(0.1488)}\approx0.366,
\]
IC del 95\%: $50.518\pm2.447(0.366) \approx (49.62,\ 51.41)$ kg.

Para el valor \textbf{individual} futuro, $\widehat{ee}\approx1.016$, y el intervalo de \textbf{predicción} del 95\% es $(48.03,\ 53.00)$ kg, claramente más ancho.
\end{frame}

\section{Resumen}

\begin{frame}{Resumen de la clase}
\begin{itemize}
\item Bajo normalidad de los errores, $\hat\beta_1\sim N(\beta_1,\sigma^2/S_{xx})$ y $\hat\beta_0\sim N\big(\beta_0,\sigma^2[\frac1n+\frac{\bar x^2}{S_{xx}}]\big)$.
\item $CME=SCE/(n-2)$ es un estimador insesgado de $\sigma^2$; con $\sigma^2$ estimado, los estadísticos estandarizados siguen $t_{n-2}$.
\item Se construyen \textbf{IC} y \textbf{pruebas $t$} para $\beta_0$ y $\beta_1$; la prueba $H_0:\beta_1=0$ es central: equivale a decir que $X$ no explica linealmente a $Y$.
\item La \textbf{prueba F} de significancia global ($CMR/CME\sim F_{1,n-2}$) es equivalente a la prueba $t$ para $\beta_1$ en regresión simple ($F=T^2$).
\item Se distinguen los IC para la \textbf{respuesta media} y el intervalo de \textbf{predicción} de una observación individual (este último siempre más ancho).
\item Próxima clase: comparación de más de dos medias mediante \textbf{análisis de la varianza}.
\end{itemize}
\end{frame}

\end{document}
